\documentclass[10pt,a4paper]{amsart}
\usepackage[margin=19mm]{geometry}
\usepackage{amsmath,amssymb,mathtools}
\usepackage[scr=rsfs,cal=euler]{mathalfa}
\usepackage{xfrac}
\usepackage[unicode,hidelinks,bookmarks=false]{hyperref}
\newcommand{\ie}{i.e.\ }
\newcommand{\cf}{cf.\ }
\newcommand{\resp}{respectively\ }
\newcommand{\Subsection}{\subsection}
\providecommand{\nobreakdash}{\nobreak\hbox{-}}
\setlength{\emergencystretch}{2em}
\multlinegap0pt
\usepackage{tikz}
\usetikzlibrary{cd}
\usepackage{mathtools}
\usepackage{braket}
\usepackage{amssymb}
\usepackage{bbm}
\usepackage[all]{xy}
\usepackage{float}
\usepackage[normalem]{ulem}
\usepackage{caption}
\usepackage{tcolorbox}
\usepackage{tikz}
\newtheorem{definition}{Definition}
\newtheorem{proposition}{Proposition}
\newtheorem{lemma}{Lemma}
\newcommand{\C}{\mathbb{C}}
\newcommand{\Hil}{\mathcal{H}}
\newcommand{\Hind}{\mathcal{H}_{\Sigma}^{\rm ind}}
\newcommand{\del}{\partial}
\newcommand{\h}{\mathcal{H}}
\title{\textbf{The Structure of the Continuum Limit of Spin Foams}}
\author{Matteo Bruno, Eugenia Colafranceschi, Fabio M. Mele, Carlo Rovelli}
\date{Source: arXiv:2603.16999v2 (CC BY 4.0)}
\begin{document}
\maketitle

\noindent\emph{Source and licence.} \textbf{The Structure of the Continuum Limit of Spin Foams}. Version arXiv:2603.16999v2 (CC BY 4.0), distributed by arXiv under the Creative Commons Attribution 4.0 International licence, http://creativecommons.org/licenses/by/4.0/, as recorded in the arXiv licence metadata for this exact version and shown on the abstract page https://arxiv.org/abs/2603.16999v2. Full licence notice: \texttt{licenses/arxiv-2603.16999-CC-BY-4.0.txt}.\par\smallskip

\noindent\textit{Editorial scope.} Counted unit: the article's proposition prop:factor (source0.tex lines 686--689). The article's complete original proof of it is delivered with it (source0.tex lines 1460--1482, one \texttt{proof} environment of 23 lines, which the article places away from the statement). No mathematics is written, repaired or re-derived: the statement and the proof are the article's own lines, in the article's own notation, and no retained line is substituted or re-flowed.  Retained uncounted as the source-local context the proof needs: lines 238--255; lines 692--704; lines 706--712; lines 1115--1118; lines 1405--1416.  Nothing was omitted from any retained span, so the package carries no omission record.  No retained line contains a citation, so the wrapper carries no bibliography and no bibliography entry.  The retained lines contain the article's own diagram code, which the wrapper typesets with the standard tikz packages; no image file is delivered and the package declares no asset dependency.  Editorial formatting only, in addition: an AMS \texttt{amsart} preamble that restates the article's own declarations and macros used by the retained lines, namely the environments \texttt{definition}, \texttt{proposition}, \texttt{lemma} on separate counters, together with the standard packages those lines need.  The retained environments print in this document as Definition 1, Proposition 1, Lemma 1 in order of appearance, whereas the article prints the same items with the numbering of its own full document.  The complete native source of the article as distributed in the arXiv e-print for this exact version is bundled unchanged as \texttt{sources/196591/source0.tex}.
\begin{definition}[Topological Quantum Field Theory]
\label{def:TQFT}
    A $d$-dimensional (unitary) TQFT is a map $Z$ which
    \begin{itemize}
        \item to any closed $(d-1)$-dimensional manifold $\Sigma$, up to orientation-preserving diffeomorphisms, assigns a Hilbert space $\mathcal H_{\Sigma}\coloneqq Z(\Sigma)$; and
        \item to the equivalence class of diffeomorphic $d$-dimensional manifolds $M$ preserving the boundary identification assigns a vector $Z(M)$ in the Hilbert space $\mathcal H_{\partial M}\coloneqq Z(\del M)$.
    \end{itemize}
    These data satisfy the following properties:
    \begin{enumerate}
        \item \textbf{Involution:} $\mathcal H_{\Sigma^*}\simeq \mathcal H_{\Sigma}^*$, where $\Sigma^*$ denotes the manifold with opposite orientation and $\mathcal H_{\Sigma}^*$ denotes the strong continuous dual space of $\mathcal H_{\Sigma}$.
        \item \textbf{Factorisation:} $\mathcal H_{\Sigma_1\sqcup\Sigma_2}\simeq \mathcal H_{\Sigma_1}\otimes\mathcal H_{\Sigma_2}$, with $\otimes$ denoting the tensor product of Hilbert spaces.
        \item \textbf{Normalisation:} When $\Sigma=\varnothing$, then $\mathcal H_{\varnothing}\simeq \C$.
        \item \textbf{Dagger:} $Z(M^*)=R_{\Sigma}(Z(M))$, where $\Sigma=\del M$ and $R_{\Sigma}:\mathcal H_{\Sigma}\to \mathcal H_{\Sigma}^*$ is the Riesz map.
        \item \textbf{Gluing:} For $M\simeq M_1\cup_{\Sigma} M_2$, with $\del M_1=\Sigma_1^*\sqcup \Sigma$ and $\del M_2=\Sigma^*\sqcup \Sigma_2$, then $Z(M)=\braket{Z(M_1),Z(M_2)}_{\Sigma}$, where $\braket{\cdot,\cdot}_\Sigma$ is the natural pairing between $\mathcal H_{\Sigma}$ and $\mathcal H_{\Sigma}^*$.
        \item \textbf{Identity:} For $M=\Sigma\times I$ with $I$ an interval, so that $\del(\Sigma\times I)=\Sigma^*\sqcup\Sigma$, then $Z(\Sigma\times I)=\mathbbm{1}_{\mathcal H_{\Sigma}}$ in $\mathcal H_{\Sigma}^*\otimes \mathcal H_{\Sigma}$.
        \item \textbf{Normalisation:} When $M=\varnothing$, then $Z(\varnothing)=1\in \C$.
    \end{enumerate}
\end{definition}

\begin{proposition}
\label{prop:dual}
    $\h_{\Sigma^*}^{\rm ind}= (\h_{\Sigma}^{\rm ind})^*$, and the following diagram
    \begin{equation}
    \label{dualcd}
    \begin{tikzcd} 
& \h_{\delta_1} \arrow{ld}[swap]{\iota_{\delta_1}} \arrow[r, "R_{\delta_1}"] &  \h_{\delta_1}^* \arrow[rd, "\iota_{\delta_1^*}"] & \\ 
      \h_{\Sigma}^{\rm ind} \arrow{rrr}[swap]{R_{\Sigma}}& & & (\h_{\Sigma}^{\rm ind})^*\\
       & \h_{\delta_2} \arrow[from=uu, crossing over, "\iota_{\delta_2\delta_1}" near start]\arrow[lu, "\iota_{\delta_2}"]\arrow{r}[swap]{R_{\delta_2}} &  \h_{\delta_2}^* \arrow{ru}[swap]{\iota_{\delta_2^*}} \arrow[from=uu, crossing over, "\iota_{\delta_2^*\delta_1^*}" near start]&  
    \end{tikzcd}
\end{equation}
commutes.~Furthermore, $\h_{\varnothing}^{\rm ind}\simeq \C$.
\end{proposition}

\begin{equation}
\label{eq:ind-net}
\begin{matrix}
    \tilde{Z}_\bullet:& \mathcal{T}_M &\to & \Hil^{\rm ind}_{\del M}\\
     & \Delta &\mapsto &\iota_{\del\Delta}\left(Z(\Delta)\right)\,,
\end{matrix}
\end{equation}

\begin{lemma}
\label{lem:contProd}
    Let $a_{\bullet}:A\to X$ and $b_{\bullet}:B\to Y$ be two nets, with $A,B$ directed posets and $X,Y$ Hausdorff topological spaces, such that $a=\lim_A a_{\bullet}$ in $X$ and $b=\lim_B b_{\bullet}$ in $Y$.~Let $f:X\times Y\to Z$ be a continuous function from $X\times Y$ to a Hausdorff topological space $Z$.~Then, $f(a,b)=\lim_{A\times B}f(a_{\bullet},b_{\bullet})$.
\end{lemma}

\begin{lemma}
    \label{lem:cutsubnet}
    Let $M=M_1\cup_{\Sigma}M_2$, with $\partial M_1=\Sigma^*_1\sqcup\Sigma,\,\del M_2=\Sigma^*\sqcup\Sigma_2$, and define
    $$
    S_{M_1,M_2,\Sigma}=\{(\Delta_1,\Delta_2)\in\mathcal{T}_{M_1}\times\mathcal{T}_{M_2}\ |\ \partial\Delta_1=\delta_1^*\sqcup\delta,\,\partial\Delta_2=\delta^*\sqcup\delta_2\  \text{for some}\ \delta_1\in\mathcal{T}_{\Sigma_1},\delta_2\in\mathcal{T}_{\Sigma_2},\delta\in\mathcal{T}_{\Sigma}\}.
    $$
    The subsposet $S_{M_1,M_2,\Sigma}$ is directed.~Furthermore,
    \begin{enumerate}
        \item the map $\gamma_{M_1,M_2,\Sigma}:S_{M_1,M_2,\Sigma}\to\mathcal{T}_{M_1\cup_{\Sigma}M_2}$ by $(\Delta_1,\Delta_2)\mapsto\Delta_1\cup_{\delta}\Delta_2$ is order-preserving and cofinal;
        \item the inclusion map $S_{M_1,M_2,\Sigma}\hookrightarrow\mathcal{T}_{M_1}\times\mathcal{T}_{M_2}$ is order-preserving and cofinal.
    \end{enumerate}
\end{lemma}

\begin{proposition}
\label{prop:factor}
    $\h_{\Sigma_1\sqcup\Sigma_2}^{\rm ind}$ and $\h_{\Sigma_1}^{\rm ind}\otimes\h_{\Sigma_2}^{\rm ind}$ are isomorphic as Hilbert spaces.
\end{proposition}

    \begin{proof}
        We have to check that the map satisfies the Atiyah's axioms as given in Def~\ref{def:TQFT}. Propositions \ref{prop:factor} and \ref{prop:dual} ensure that $\Sigma\mapsto\Hind$ satisfies Involution, Factorisation, and Normalisation. Hence, we just need to verify the properties of the limit $Z(M)$.
        Normalisation is immediate:
        \[Z(\varnothing)=\lim_{\Delta\in\mathcal{T}_{\varnothing}}\tilde{Z}_{\Delta}=\tilde{Z}_{\varnothing}=\iota_{\varnothing}(Z(\varnothing))=1.\]
        The Dagger property follows from the continuity of the Riesz map and the Dagger property of $\tilde{Z}$, namely
        \[Z(M^*)=\lim_{\Delta\in\mathcal{T}_{M^*}}\tilde{Z}_{\Delta}=\lim_{\Delta\in\mathcal{T}_{M}}\tilde{Z}_{\Delta^*}=\lim_{\Delta\in\mathcal{T}_{M}}R_{\Sigma}(\tilde{Z}_{\Delta})=R_{\Sigma}(\lim_{\Delta\in\mathcal{T}_{M}}\tilde{Z}_{\Delta})=R_{\Sigma}(Z(M)).\]
        The gluing property requires a bit more discussion. Consider $M=M_1\cup_{\Sigma}M_2$, with $\del M_1=\Sigma_1^*\sqcup\Sigma$ and $\del M_2=\Sigma^*\sqcup\Sigma_2$, then
        \begin{align*}
    Z(M)=\lim_{\Delta\in\mathcal{T}_M}\tilde{Z}_{\Delta}=\lim_{(\Delta_1,\Delta_2)\in S_{M_1,M_2,\Sigma}}\tilde{Z}_{\Delta_1\cup_{\delta}\Delta_2}=\lim_{(\Delta_1,\Delta_2)\in S_{M_1,M_2,\Sigma}}\braket{\tilde{Z}_{\Delta_1},\tilde{Z}_{\Delta_2}}_{\Sigma}\;,
\end{align*}
where, in the second equality, we used Lemma~\ref{lem:cutsubnet} together with the equality of the limits of a net and of all of its subnets; while, in the last equality, we used the gluing property of $\tilde Z_{\bullet}$ from the main text (cfr.~property 3 below Eq.~\eqref{eq:ind-net}).
    Let us now look at the product $\braket{Z(M_1),Z(M_2)}_{\Sigma}$. Using Lemma~\ref{lem:contProd}, with the continuous function given the natural pairing $\braket{\cdot,\cdot}_{\Sigma}$, and Lemma~\ref{lem:cutsubnet}, we have
    \begin{align*}
\braket{Z(M_1),Z(M_2)}_{\Sigma}&=\left\langle\lim_{\Delta_1\in\mathcal{T}_{M_1}}\tilde{Z}_{\Delta_1},\lim_{\Delta_2\in\mathcal{T}_{M_2}}\tilde{Z}_{\Delta_2}\right\rangle_{\Sigma}\\
    &=\lim_{(\Delta_1,\Delta_2)\in\mathcal{T}_{M_1}\times\mathcal{T}_{M_2}}\braket{\tilde{Z}_{\Delta_1},\tilde{Z}_{\Delta_2}}_{\Sigma}\\
    &=\lim_{(\Delta_1,\Delta_2)\in S_{M_1,M_2,\Sigma}}\braket{\tilde{Z}_{\Delta_1},\tilde{Z}_{\Delta_2}}_{\Sigma}
\end{align*}
Then, from the unicity of the limit, we conclude that
\[Z(M_1\cup_{\Sigma}M_2)=\braket{Z(M_1),Z(M_2)}_{\Sigma}.\]
This holds true also for the cylinder $M_{1,2}=\Sigma\times I$, from which it follows that $Z(\Sigma\times I)$ is a projector: 
\[Z(\Sigma\times I)=\braket{Z(\Sigma\times I),Z(\Sigma\times I)}_{\Sigma}.\]
Thus, the theory specified by the map under consideration defines a semi-functor for $\mathsf{Cob}_{\rm d}^{\rm or}$ to $\mathsf{Hilb}$. Turning the semi-functor into functor by the standard procedure, we can just consider as Hilbert space $\mathrm{Ran}(Z(\Sigma\times I))$ and get also the Identity property of a TQFT.
\end{proof}

\end{document}
